Dieses Experiment untersucht die Lebensdauer des $\frac{5}{2}^+$-Zustands von $^{133}_{55}$Cs. Zwei Szintillationsdetektoren werden hierfür eingesetzt, deren Signale digital und mit einem Oszilloskop analysiert werden. Zunächst wird die zeitliche Auflösung des Detektionssystems durch die Aufzeichnung der Koinzidenzphotonen aus dem Zerfall von $^{22}_{11}$Na bestimmt. Anschließend wird die Zeitdifferenz zwischen den beiden wahrscheinlichsten Zerfällen von $^{133}_{56}$Ba gemessen, woraus sich die Lebensdauer ableiten lässt. Abschließend wird die Energieauflösung des Messgeräts bestimmt. 